\section{Реестр решённых и нерешённых вопросов}
В этой работе важно различать пять статусов: исходный сертифицированный
результат репозитория; его доказанное следствие; новое утверждение рабочего
аппендикса с доказательством; конечный вычислительный результат; открытый
вопрос. Неудача ограниченного поиска не считается доказательством
невыпрямимости. Несколько эквивалентных формулировок одного кандидата
не считаются решениями нескольких общих проблем.

\begin{longtable}{@{}L{0.24\linewidth}L{0.40\linewidth}L{0.30\linewidth}@{}}
\toprule Вопрос & Точный итог & Граница вывода\\\midrule\endhead
Трилистник Шастри & Предъявлены взаимно обратные выпрямители над
$K=\Q(\sqrt{-3})$. & Опровергается невыпрямимость именно этой комплексной
модели; вещественный узел остаётся нетривиальным.\\[5pt]
Гармоническая восьмёрка & Выпрямлено $H(3,5,7)$ с линейным коэффициентом
$+5u$ во второй компоненте. & Эта модель отличается от Example 2
Шастри: там средняя компонента равна $u^5-5u^3+4u$.\\[5pt]
Conjecture 5.1 в \cite{vdER} & Для указанного там $D$ ядро изоморфно
$\C^{[4]}$, в противоположность гипотезе. & Конкретный кандидат не является
контрпримером к сокращению.\\[5pt]
Conjecture 5.2 в \cite{vdER} & Исходный срез является координатой;
его нулевой слой есть $\A^4_\C$. & Следствие полиномиальности того же ядра.\\[5pt]
Conjecture 5.3 в \cite{vdER} & Связанное расширенное кольцо Риса
полиномиально. & В статье это конкретная гипотеза Асанумы;
другие гипотезы Асанумы не затрагиваются.\\[5pt]
Гармонические семейства & Все допустимые $H(3,4,c)$ и $H(3,5,c)$
выпрямляются над $K$. & Не классификация всех $H(a,b,c)$ и не утверждение
о бесконечном числе новых типов вещественных узлов.\\[5pt]
Степени срезов & Полное конечное линейное сведение, восемь точных минимумов;
вся зависимость от кратности для двух моделей. & Степень зависит от исходных
координат и не является инвариантом сопряжения.\\[5pt]
Нормальное поле & Доказан критерий выпрямления через локальную конечность
и конструктивное восстановление координат. & Существование такого поля
для произвольного вложения ещё не доказано.\\[5pt]
$H(3,7,8)$ и $H(4,5,7)$ & Сертифицированы пространства полей-кандидатов,
барьеры степеней и несколько исключённых классов. & Комплексная выпрямимость,
полиномиальность новых ядер и их полный ML здесь не установлены.\\[5pt]
Общие задачи & Пространственная проблема вложения, все гармонические
тройки, сокращение и общая линеаризация здесь не решены. & Утверждение
«не решено этой работой» не заменяет отдельного обзора всех известных
вариантов этих проблем.\\\bottomrule
\end{longtable}

\paragraph{Историческая точность.}
Плоская теорема Абъянкара--Моха--Судзуки уже является классическим
результатом \cite{AM,Suzuki}. Историческое предположение о существовании
невыпрямляемых пространственных прямых не опровергается выпрямлением
одного или двух кандидатов. Вопросы 1 и 2 статьи Шастри \cite{Shastri}
о существовании полиномиальных представлений узлов с комплексным
вложением решены самим Шастри. Наш вывод относится к конкретному
предположению о его трилистнике, зафиксированному в \cite[с.~1319]{vdER}.

\paragraph{Не заявляемые результаты.}
Не доказаны общая гипотеза Абъянкара--Сатайе, прирученность либо дикость
построенных автоморфизмов, минимальность степеней самих выпрямителей,
общая нижняя граница $6n+1$ для степеней срезов и точная общая сложность
разложения последовательных модификаций. Эти вопросы сохранены в
исследовательской программе, но не включены в список закрытых.
